\abstract
Dependency-Injection is a technique used to add and load dependencies within an application dynamically.
Being able to alter an application dynamically can be both, wanted and unwanted.
This research report will evaluate if enabling Dependency-Injection in Java poses higher risks and threads or if the benefits are worth taking risks.
Additionally, it might be possible to remove the risks and threads altogether.
It should be evident that alerting a running application, especially dynamically, can be of high-security risk.

Initially, we planned to conduct a qualitative empirical experiment, showcasing different levels of security, and evaluate based on these results what measures can be taken to reduce the overall risk.
However, later we decided that informing about the potential risks, threats, attacks and exploits is a better approach for this report.
Doing so, we do not prove how one can inject malicious code, but what is possible when dynamically loading and injecting code into an application.
Additionally, different types of injecting (``attacking'') will be discussed.

We will first cover what \textit{Dependency-Injection} is and define some terms we will use in the report.
During that, we talk about different scenarios where \textit{Dependency-Injection} applies.
Next, we explain, on a basic level, how one can implement \textit{Dependency-Injection} in \textbf{Java}, how to handle it and conduct our experiment.
Afterwards, we present and evaluate our results and finally discuss what these results mean.

In summary, the report concludes that dependency-injection is indeed possible and, without further securing an application, can pose a high risk not only to the application, but also the system and user.
However, there are also benefits, such as being able to extend an application to the users' needs dynamically.

\clearpage